\section{Introduction}

In the field of Natural Language Processing (NLP), plain language summarization (PLS) distills complex documents, such as scientific articles, into accessible summaries for non-expert audiences while preserving essential meaning~\cite{chandrasekaran2020overview}. The COVID-19 pandemic highlighted the critical need to make scientific knowledge accessible to the general public~\cite{wang-etal-2020-cord}. By enhancing public engagement with research, PLS can help bridge the gap between expert knowledge and general understanding.


Although human evaluation remains the gold standard for assessing summary quality and readability, the high cost and slow turnaround~\cite{Liu2022RevisitingTG} have led many researchers to rely on automatic evaluation metrics for evaluating PLS summaries~\cite{goldsack2022making,Guo2020AutomatedLL}. Although these metrics have been validated in fields such as education and law~\cite{gray1935makes, Han2024TheUO}, their effectiveness in reflecting readability in the context of PLS remains unproven.


Are automated readability metrics appropriate evaluators for the task of PLS? We explore whether the definition of readability as implemented by automated measures matches the definition used by the PLS research community. Additionally, given that Language Models (LMs) can reason over complex language tasks~\cite{Brown2020LanguageMA, Wei2022ChainOT, Yang2024DoLL}, we explore whether LMs can judge the readability of a summary. Given these motivations, we ask the following research questions (RQs).

\paragraph{{\rqone\label{rq1}}What is the current standard of evaluation in PLS literature?} 
We review PLS literature by collecting relevant papers published in *ACL venues from 2013 to 2025 and note the readability evaluation method used in the study. We find that the majority of papers focus on a small number of traditional readability metrics, such as Flesch-Kincaid grade Level (FKGL)~\cite{flesch1952simplification}.
This finding motivates our analysis of the suitability of traditional readability metrics for PLS evaluation.


\paragraph{{\rqtwo\label{rq2}}How well do traditional readability metrics correlate with human readability judgments?} Since the PLS research community employs these traditional metrics (RQ1), we assess their suitability by measuring their correlation with human readability judgments. A low correlation would suggest that a metric is inadequate for evaluating PLS readability, and would necessitate the PLS research community identify and move to better metrics. To the best of our knowledge, this work is the first to compare traditional readability metrics to human readability judgments for PLS.

\paragraph{{\rqthree\label{rq3}}How well do LM-based evaluators correlate with human readability judgments?} Traditional readability metrics primarily use lexical features, such as the number of syllables in a word, to measure readability. In contrast, LMs may capture more complex attributes of readability than traditional metrics, such as the inclusion of necessary context and explanation of key concepts. The findings of this research question have important implications for both the best practices in evaluation of PLS and the broader NLP community's understanding of LM capabilities.

\paragraph{{\rqfour\label{rq4}} What do LM-based evaluators reveal about the readability of popular summarization datasets?} Researchers often rely on traditional readability metrics when assessing summaries in new methods or datasets. However, if these metrics correlate poorly with human judgments, the resulting conclusions may be flawed. Similarly, existing datasets, which often arise from data of convenience, may be poorly suited to PLS research. This RQ explores what LM-based evaluators reveal about the readability of popular summarization datasets and how LM-based conclusions differ from those based on traditional readability metrics. 

We answer these questions through the following contributions. First, we survey PLS papers published in *ACL venues and find that the most popular metric for readability evaluation is Flesch-Kincaid Grade Level (FKGL)~\cite{flesch1952simplification}. Motivated by these findings, we then compare 8 traditional readability metrics to human judgments. We show that 6 of the 8 metrics have a poor correlation (less than 0.3 Pearson correlation) with human judgments, including FKGL, indicating that these metrics are poor measures of readability for PLS. Additionally, we compare the judgments of 5 LMs to human judgments and show that LMs outperform the traditional metrics. We demonstrate that LMs have promising potential as evaluators by reasoning over more complex attributes of readability. We use LM evaluators to re-evaluate 10 summarization datasets and show that some summarization datasets intended for PLS achieve similar readability scores to those aimed at expert audiences, calling into question the utility of these data. Finally, based on a thorough analysis of current readability evaluation practices, we offer recommendations for best practices in PLS evaluation and identify opportunities for future work. 


	% 5.	Expand the Introduction for a Slower, More Detailed Development
	% •	The introduction should be longer, ideally extending into a second column on the next page, to ensure a thorough setup of the research problem and motivation.
	% 6.	Explicitly Write Out Research Questions with Explanations
	% •	Each research question should be clearly stated, followed by a brief explanation:
	% •	RQ1: How well do traditional readability metrics correlate with human judgments?
	% •	If they do not correlate, then they are not suited for evaluating the success of plain language summarization.
	% 7.	Clarify the Description of Tasks
	% •	When referring to “the nine tasks,” provide a brief explanation of their nature. For example:
	% •	These tasks span various scientific domains, with content aimed at different audiences, such as children or journalists

